\section{构建 AI Agent 的关键模块}

\subsection{Vertex AI 平台架构}

Google Vertex AI 的设计体现了企业 AI 平台的共性需求：

\begin{itemize}
    \item \textbf{芯片层}：GPU/TPU 选择灵活性（当前全球芯片严重短缺）
    \item \textbf{模型层}：第一方 + 第三方 + 开源模型的统一接口
    \item \textbf{工具层}：Model Builder 和 Agent Builder 提供定制和编排能力
    \item \textbf{应用层}：搜索、对话、数据分析等垂直应用
\end{itemize}

\subsection{Gemini 的差异化}

\begin{itemize}
    \item 原生多模态（从训练开始就混合多种模态）
    \item 超大上下文窗口（实验中已达 1000 万 token）
    \item Needle-in-a-haystack 测试：在超长上下文中精确检索信息
\end{itemize}

\subsection{Planner 与 Memory 模块}

Agent 内部通常包含 planner、executor、memory 三个协同组件。典型流程是：

\begin{itemize}
    \item \textbf{Planner}：把用户意图拆解为步骤（比如搜索、工具调用、总结）
    \item \textbf{Memory}：记录每次对话、查询与工具输出，提供 context re-use
    \item \textbf{Executor}：负责调用特定工具、API 或自定义运行时
\end{itemize}

\begin{importantbox}{planner+memory 的同步挑战}
当 planner 制定步骤后，memory 需要快速返回过去的执行结果。如果 memory 不够精细（只保留最近一轮），planner 会重复调用工具；如果 memory 太大，又会造成 latency。解决办法是引入 compressed memory + recall hints，即在 memory 中维护 query summary 供 planner 参考。
\end{importantbox}
\begin{importantbox}{Planner/Memory 的分工}
Memory 为 planner 提供上下文缓存，planner 决策要素（上一轮失败、工具成功率），executor 在安全 guardrail 之下运行。这个三层架构让 agent 能在长上下文、跨工具的任务中保持连贯。
\end{importantbox}

\subsection{Tool Orchestration 与 Workflow}

Agent 要协调多个工具（搜索、数据库、API），需要明确工作流和 tool specification：

\begin{itemize}
    \item 定义 tool intent、inputs、outputs，避免 model hallucinating tool usage
    \item 用 `tool descriptor` 验证 LLM 的调用参数（类型、范围）
    \item 在 workflow 中插入 `pre-conditions` 与 `post-conditions`
\end{itemize}

\begin{knowledgebox}{workflow 的 state machine 设计}
Complex agent often models workflow as state machine：每个状态代表一个任务阶段（理解、planning、execution、confirmation），transition 触发条件来自 planner 的判断。这样可以在开发阶段用 unit test 模拟不同分支，提高 system reliability。
\end{knowledgebox}
\begin{knowledgebox}{tool descriptor 的价值}
当模型需要调用 `fetch\_legislation()` 或 `run\_cad\_simulation()` 时，tool descriptor 提供 schema、authority、expected latency，避免 agent 盲目调用或超时。
\end{knowledgebox}

\subsection{本章小结}

企业 AI 平台需要在芯片、模型、工具和应用各层提供选择灵活性。

